%=========================================
% Section: Literature Survey
% All citations linked to references.bib
%=========================================

\section{Literature Survey}

This section reviews influential studies relevant to electric vehicle charging
station (EVCS) planning, multi-objective evolutionary algorithms, queueing-based
service modelling, stochastic demand forecasting, and comparative algorithm
evaluation, arranged in reverse chronological order.

\begin{enumerate}

% -------------------------------------------------------------------------
\item \textbf{Palka et al.\ \cite{palka_planning_2025}} develop an integrated
EVCS scheduling and planning system combining machine learning demand prediction
with Mixed-Integer Linear Programming (MILP) for station sizing and placement.
Their hybrid AI--optimisation framework reduces planning cost by $\sim$12\%
over purely heuristic approaches. The formulation remains single-period and
single-objective, motivating the evolutionary four-objective framework of this
thesis.

% -------------------------------------------------------------------------
\item \textbf{Wang et al.\ \cite{wang_deadline_2025}} propose an online
deadline-aware EV charging model accounting for state-of-charge (SoC) dependent
peak charging rates, showing that ignoring SoC dynamics underestimates peak grid
load by up to 18\% under fluctuating demand. Their load-stability analysis
reinforces the need for finite-buffer M/M/$c$/K constraints across Time-of-Use
periods adopted in this thesis.

% -------------------------------------------------------------------------
\item \textbf{Zhao et al.\ \cite{zhao_optimal_2023}} propose a bi-level EVCS
location model for self-charging and valet-charging user modes. Their
Monte~Carlo demand integration shows ignoring uncertainty yields over-optimistic
coverage estimates, directly informing the 100-scenario stochastic simulation
module used here. Pricing, grid capacity, and queueing constraints are absent.

% -------------------------------------------------------------------------
\item \textbf{Li et al.\ \cite{li_joint_2022}} jointly optimise charging
station placement and pricing under stochastic arrivals via chance-constrained
programming, highlighting pricing as an integrated objective. This motivates
replacing raw profit with net ROI (annual profit / CAPEX) to avoid Pareto
front degeneracy caused by near-perfect linear correlation with coverage.

% -------------------------------------------------------------------------
\item \textbf{Xia et al.\ \cite{xia_grid_2022}} show that ignoring power
network constraints in EVCS placement yields infeasible solutions at high EV
penetration, and penalise solutions exceeding zone-level substation capacity.
This per-zone grid capacity constraint is adopted explicitly in the problem
formulation of this thesis.

% -------------------------------------------------------------------------
\item \textbf{Kulkarni et al.\ \cite{kulkarni_bass_2021}} apply the
generalised Bass logistic diffusion model to forecast EV adoption in India,
estimating an inflection point near 2030 consistent with NITI~Aayog targets.
Their parameterisation ($t_0 = 2030$, $a = 0.3$) is adopted in the logistic
adoption curve and validated against ward-level demand projections from
Census~2011 data.

% -------------------------------------------------------------------------
\item \textbf{Paul and Bhattacharya \cite{paul_nsga2_2021}} study
multi-objective EVCS siting for Kolkata using NSGA-II with cost and coverage
objectives, confirming evolutionary superiority over classical programming for
large-scale Indian city instances. The work is restricted to two objectives
with no queueing, ROI, or stochastic demand --- gaps directly addressed here.

% -------------------------------------------------------------------------
\item \textbf{Bernardo et al.\ \cite{bernardo_density_2021}} demonstrate that
urban density significantly affects appropriate service radius: dense centres
require $\sim$1--2~km versus $\sim$3--5~km in suburban areas. This informs the
zone-adjusted radii of 1.5~km (Central Core) and 2.5~km (Outer zones) adopted
following MoRTH guidelines.

% -------------------------------------------------------------------------
\item \textbf{Zhang et al.\ \cite{zhang_mmck_2019}} show that ignoring finite
system capacity underestimates blocking probability by 15--20\% at high
utilisation, justifying the M/M/$c$/K model with $K = c + W$ waiting bays over
the simpler Erlang-B loss model. Their capacity planning framework forms the
analytical basis of the queueing evaluation layer in this thesis.

% -------------------------------------------------------------------------
\item \textbf{Tan et al.\ \cite{tan_tou_2016}} introduce a three-period
Time-of-Use (ToU) arrival rate model --- Peak, Normal, and Idle --- accurately
capturing the bimodal daily demand pattern in EV charging field data. This
thesis extends their framework by embedding ToU-weighted M/M/$c$/K evaluation
directly inside the evolutionary loop rather than as post-hoc analysis.

% -------------------------------------------------------------------------
\item \textbf{Sadeghi-Barzani et al.\ \cite{sadeghi_fast_2014}} optimise
placement and sizing of fast charging stations by minimising installation cost
and power loss under grid-side constraints. The work is single-objective and
ignores user accessibility, queue performance, and financial viability ---
limitations that motivate the multi-objective formulation herein.

% -------------------------------------------------------------------------
\item \textbf{Reyes-Sierra and Coello~Coello \cite{reyes_sierra_2006}} survey
NSGA-II versus MOPSO across multiple problem domains, finding MOPSO converges
faster early but sacrifices spread diversity, while NSGA-II maintains a
better-distributed Pareto front. Their theoretical prediction is verified in
this thesis through formal HV, GD, and Spread evaluation on the binary EVCS
placement instance.

% -------------------------------------------------------------------------
\item \textbf{Coello~Coello and Lechuga \cite{coello_mopso_2002}} propose
MOPSO, extending Particle Swarm Optimisation to multiple objectives via an
external non-dominated archive and crowding-based leader selection. Its
weakness on binary combinatorial problems --- where continuous velocity must
be binarised via sigmoid threshold --- is examined in the results chapter.

% -------------------------------------------------------------------------
\item \textbf{Deb et al.\ \cite{deb_nsga2_2002}} introduce NSGA-II with fast
non-dominated sorting ($\mathcal{O}(MN^2)$) and crowding distance diversity
preservation. Its elitism, convergence guarantees, and compatibility with
binary decision variables make it the primary optimiser in this thesis.
Queueing and pricing constraints are not natively incorporated within the
evaluation loop.

\end{enumerate}

%------------------------------------------------------------------
\section{Research Gaps}
\label{sec:lit_gaps}

The review of~\cite{palka_planning_2025,wang_deadline_2025,zhao_optimal_2023,
li_joint_2022,xia_grid_2022,kulkarni_bass_2021,paul_nsga2_2021,
bernardo_density_2021,zhang_mmck_2019,tan_tou_2016,sadeghi_fast_2014,
reyes_sierra_2006,coello_mopso_2002,deb_nsga2_2002} reveals five
methodological gaps that collectively motivate the proposed framework:

\begin{enumerate}
  \item \textbf{Indian city focus with real ward-level data.} Most
        multi-objective EVCS studies~\cite{zhao_optimal_2023,li_joint_2022,
        xia_grid_2022} use Chinese or European city models. Studies for
        Tier-1 Indian cities with actual ward-level demographic data and
        India-specific policy parameters are scarce; \citet{paul_nsga2_2021}
        is a partial exception but is restricted to two objectives.

  \item \textbf{Integrated four-objective formulation.} Existing studies
        address two~\cite{paul_nsga2_2021} or three~\cite{zhao_optimal_2023}
        objectives. The simultaneous optimisation of capital expenditure,
        spatial coverage, net ROI, and queue waiting time in a single Pareto
        framework has not been reported for Indian cities.

  \item \textbf{M/M/$c$/K with ToU pricing embedded in the optimization loop.}
        Prior work either uses over-simplified Erlang-B models or applies
        time-averaged M/M/$c$ queues post-hoc~\cite{tan_tou_2016}. An
        M/M/$c$/K model~\cite{zhang_mmck_2019} with ToU demand periods and a
        hard blocking probability constraint ($P_{\text{block}} \leq 10\%$)
        embedded directly inside the evolutionary evaluation represents a
        methodological advance.

  \item \textbf{NSGA-II vs.\ MOPSO formal comparison on binary placement
        problems.} While comparative surveys exist~\cite{reyes_sierra_2006},
        formal empirical comparison of NSGA-II~\cite{deb_nsga2_2002} and
        MOPSO~\cite{coello_mopso_2002} using HV, GD, and Spread on a binary
        EVCS placement instance with queueing constraints has not been
        reported.

  \item \textbf{Net ROI as an independent objective.} As highlighted by
        \citet{li_joint_2022}, using raw profit creates near-perfect linear
        correlation with coverage ($r > 0.99$), yielding a degenerate Pareto
        front. Replacing it with net ROI (annual profit / CAPEX) creates
        genuine trade-off diversity in the multi-objective search space.
\end{enumerate}

%------------------------------------------------------------------
\section{Comparative Summary of Key Literature}

Table~\ref{tab:lit_survey} provides a structured comparison of all fourteen
reviewed studies across methodology and key contributions.

\begin{table}[htbp]
\centering
\small
\setlength{\tabcolsep}{4pt}
\renewcommand{\arraystretch}{1.3}
\begin{tabular}{|c|p{9em}|p{10em}|p{17em}|}
\hline
\textbf{No.} & \textbf{Study} & \textbf{Method / Model} &
\textbf{Key Contribution / Limitation} \\
\hline
1  & \cite{palka_planning_2025}
   & ML + MILP hybrid
   & Cost-efficient scheduling; single-period, single-objective. \\
\hline
2  & \cite{wang_deadline_2025}
   & Online SoC-aware model
   & Grid-stability under fluctuating load; no spatial siting or Pareto
     optimisation. \\
\hline
3  & \cite{zhao_optimal_2023}
   & Bi-level + Monte~Carlo
   & Robust demand modelling; no pricing or queueing. \\
\hline
4  & \cite{li_joint_2022}
   & Joint placement + pricing
   & Stochastic pricing; no evolutionary multi-objective. \\
\hline
5  & \cite{xia_grid_2022}
   & Grid-constrained placement
   & Power-feasible siting; single-objective. \\
\hline
6  & \cite{kulkarni_bass_2021}
   & Bass logistic adoption, India
   & India-calibrated $t_0=2030$; no spatial model. \\
\hline
7  & \cite{paul_nsga2_2021}
   & NSGA-II, Indian city
   & Two-objective only; no queueing or ROI. \\
\hline
8  & \cite{bernardo_density_2021}
   & Density-based radius model
   & Zone-adjusted radii; no optimization framework. \\
\hline
9  & \cite{zhang_mmck_2019}
   & M/M/$c$/K capacity planning
   & Finite-buffer blocking; no spatial siting. \\
\hline
10 & \cite{tan_tou_2016}
   & ToU queueing model
   & Time-of-use arrivals; not embedded in optimization. \\
\hline
11 & \cite{sadeghi_fast_2014}
   & Fast-charging siting + grid
   & Grid-aware siting; single-objective only. \\
\hline
12 & \cite{reyes_sierra_2006}
   & NSGA-II vs.\ MOPSO survey
   & NSGA-II superior spread; no placement application. \\
\hline
13 & \cite{coello_mopso_2002}
   & MOPSO
   & Fast convergence; weaker spread on binary problems. \\
\hline
14 & \cite{deb_nsga2_2002}
   & NSGA-II
   & Benchmark MOEA; no queueing or pricing integration. \\
\hline
\end{tabular}
\caption{Comparative summary of key studies arranged by year (latest first) on
EVCS siting, queueing, multi-objective optimization, and demand forecasting.
Full bibliographic details in the reference list.}
\label{tab:lit_survey}
\end{table}